\chapter{Validation, Quality, Discussion, and Limitations}
\label{chap:validation}

\section{Validation scope and evidence date}

The software-validation evidence in this chapter begins with the full application audit performed on 31 July 2026, includes the foundation-model work completed on 14 August and the coordinated Azure promotion completed on 15 August, and concludes with repository-wide remediation and local revalidation on 16 August 2026. Historical reports are retained for chronology because test counts, routes, and capabilities changed over time. The audit uses backend tests, frontend quality commands, browser tests, artifact checks, real-wrapper smoke tests, source review, immutable image digests, and public cloud readiness. GitHub Actions run \texttt{31893388429} verifies source commit \texttt{d25295b}; Azure revisions with the same suffix serve that verified web, API, and email-worker release. The later local remediation is not described as deployed until a new remote promotion is performed.

\section{Backend validation}

The SQLite repository run passed 142 tests with four expected database-specific skips and measured 82\% application coverage. The new cases cover the month manifest and checksum, 365-day input preparation, 30-target inference, cadence-aware API persistence and reload, monthly PDF output, hourly challenger protocol, and application security headers. A focused real-wrapper smoke test also loaded the exact Chronos-2 base and adapter and returned 30 finite ordered quantiles. A uniquely named disposable PostgreSQL 16 database was then migrated through all 18 Alembic revisions; the full suite passed 146 tests with no skips and \texttt{alembic check} reported no schema drift. The new ownership-finalization revision refuses unresolved orphan readings before enforcing a non-null meter reference.

The PostgreSQL run is the stronger result because it exercises locking and database behavior that SQLite cannot reproduce. It also confirms that migrations can construct the schema from empty state. The audit did not mutate the live product database.

Coverage is reported only for the measured application package and is enforced in CI with a 70\% minimum. The observed 82\% does not imply that every path, provider, or operational failure mode is exercised.

\section{Frontend and browser validation}

ESLint passed with hooks, immutability, unused-variable, and explicit-any rules enforced as errors; TypeScript type checking passed; and the production build completed with 25 routes. The current Playwright matrix executes 204 tests across desktop Chromium, Firefox, WebKit, and three responsive projects. It covers 360-pixel, 390-pixel iPhone, and 768-pixel profiles. Simulator coverage verifies the one-year hourly bootstrap, explicit synthetic provenance, all-model readiness message, reset isolation, and handoff to real CSV data. Settings coverage additionally verifies that a zero budget is preserved and that language changes persist. Populated 167-hour and 169-hour local-week tests continue to exercise daylight-saving boundaries.

\begin{table}[H]
\centering
\caption{Current software validation summary.}
\label{tab:testsummary}
\begin{tabularx}{\textwidth}{p{0.27\textwidth}p{0.25\textwidth}Y}
\toprule
Area & Result & Interpretation \\
\midrule
Backend and training, SQLite & 142 passed, 4 skipped; 82\% coverage & PostgreSQL-only behavior excluded; CI floor is 70\% \\
Backend and training, PostgreSQL & 146 passed, 0 skipped & Empty database migrated through 18 revisions; zero schema drift \\
ML-enabled backend test image & 117 passed, 4 skipped & Hash-locked foundation stack; backend-only container scope \\
Frontend lint and type checking & Passed & Static quality gates are green \\
Frontend production build & Passed, 25 routes & Current route tree compiles for production \\
Playwright six-project matrix & 204 passed, 0 failed & Desktop and three responsive profiles green \\
Docker Compose & Seven-service configuration valid & Current source topology includes a dedicated simulation worker \\
Azure public deployment & Web/API/worker \texttt{v4-d25295b}; healthy & Managed PostgreSQL and all three model horizons warmed; simulation/alert/avatar workers remain absent \\
\bottomrule
\end{tabularx}
\end{table}

\section{Model-contract and artifact validation}

The product does not rely only on a database model name. Artifact tests verify manifest parsing, supported capability, checksum, strict load, CPU inference, finite output, expected shape, cadence, and reference points where recorded. The two active TFT checkpoints share a file size of 5,600,105 bytes but have distinct SHA-256 hashes. The month release separately verifies a 4,854,584-byte LoRA adapter, the pinned Chronos-2 base revision, FP32 inference, and exactly 30 ordered daily quantiles. A real-wrapper smoke test completed a checksum-verified load and prediction.

The manifests also preserve limitations. London-cohort metrics are not a guarantee for a client site. New hourly requests use their recent 336-hour window for normalization, while month requests require a substantially longer daily history. The month model's positive Tetouan and Portugal evidence coexists with a negative southern-Morocco diagnostic, and its conformal correction is not a per-site guarantee. The hourly Chronos adapters pass their London gates but remain inactive because deployment value includes memory, startup, input migration, and geographic evidence as well as MAE.

\section{Security validation}

The audit reviewed password hashing, JWT and refresh-token behavior, account-action tokens, Google identity gating, role checks, ownership filters, upload bounds, request validation, SSRF controls, and secret configuration. CI includes a full-history gitleaks job. No P0 authentication bypass or data-loss defect was identified. Permanent account deletion is additionally protected by recent password reauthentication and a literal confirmation value. Its implementation removes the complete ownership graph with ordered set-based statements, then writes only an anonymized deletion audit event, avoiding the request timeout observed with the earlier row-by-row ORM cascade.

Both application layers now define response security policy. Next.js emits CSP, HSTS, frame denial, MIME-sniffing prevention, referrer policy, permissions policy, and cross-origin opener policy; FastAPI emits the applicable API-side headers. Browser and API smoke tests verify the policy, and expensive forecast, simulation, import, export, and deletion operations have explicit rate limits. Interactive OpenAPI and ReDoc discovery are available in debug mode and disabled by production policy. New, reset, and changed passwords share a 12-character minimum; existing hashes remain valid for login. The final npm audit and the hash-locked runtime, development, and foundation Python audits reported no known vulnerabilities; the custom CPU Torch wheel is outside PyPI's advisory database and is identified separately rather than silently treated as audited. A full-history scan covered 407 commits and found no leaked secrets after review of exact documentation false positives. Google integration uses only the public client ID; no provider secret is required or stored by EnergyAI.

\section{Docker and deployed service health}

The current seven-service Compose configuration resolves successfully and separates simulator advancement from the API lifecycle. The final backend test image installed every production and foundation dependency through hash-locked files, prefetched the exact base revision, passed its internal backend suite, ran as UID 10001, and contained no compiler. The standalone frontend image passed the same non-root check and returned its security headers in a live HTTP smoke test. GitHub Actions and Azure evidence below refer to the preceding promoted release; the remediated source requires a new image promotion and explicit simulation-worker provisioning before those local results become cloud evidence.

The final Azure smoke test for the preceding release confirmed public TLS and HTTP 200 for the principal public pages, API liveness and readiness, managed-database connectivity, warmed 24-hour, 168-hour, and 720-hour models, ready SMTP capability, and a healthy one-replica email worker. Google identity remains published for external users; the earlier account-chooser test created no user. A read-only Azure check found \texttt{ca-energyai-web--v4-d25295b}, \texttt{ca-energyai-api--v4-d25295b}, and \texttt{ca-energyai-email-worker--v4-d25295b} healthy, with 100\% traffic on the web and API release. The web and API retain zero-to-one scaling; the email worker remains at one replica. Operational readiness is therefore demonstrated for that controlled public PFE presentation, while promotion of the remediated images, simulation/alert/avatar workers, durable avatar storage, centralized monitoring, restore evidence, load testing, and disaster-recovery procedures remain conditional.

\section{Requirement-to-test traceability}

The report's traceability registry maps each major requirement to code and validation evidence. Examples include JWT authentication to authentication tests, RBAC to admin authorization tests, per-user data isolation to PostgreSQL integration tests, CSV preview/import and simulator handoff to ingestion/browser tests, persistent demonstration continuity to simulator tests, bounded account deletion to ownership-graph tests, and each forecast horizon to model-contract and dual-horizon browser tests. The complete machine-readable matrix is stored in \evidence{report/evidence/requirements\_traceability.csv}.

Traceability does not remove the need for judgment. For example, a route and test can show that account deletion works, but legal compliance requires institutional and jurisdiction-specific review. Likewise, an artifact test proves file identity and output shape, not real-world forecast calibration.

\section{Scientific limitations}

\subsection{Geographic transfer}

The final model evidence comes from Low Carbon London. Climate, building characteristics, tariffs, appliance ownership, and daily routines can differ materially in Morocco. Cold-start evaluation within London is stronger than random household windows, but it is not cross-country validation. A Moroccan pilot cohort is required before making local performance guarantees.

\subsection{Input scope}

The product contract uses timestamp and energy. This lowers onboarding friction but excludes weather, occupancy, tariff events, and building metadata that may explain regime changes. Calendar features are deterministic, not measurements. Future work should compare the marginal value of weather forecasts and static site attributes under a strictly causal protocol.

\subsection{Uncertainty calibration}

The TFT produces 0.1, 0.5, and 0.9 quantiles, but the manifests state that they are not recalibrated for each client. Reanalysis under the original research scaler finds 80.7963\% empirical coverage for the nominal central 80\% interval at 24 hours and 78.3518\% at 168 hours, with mean widths of 0.559770 and 0.585659 kWh. Under serving normalization, coverage is 81.0108\% and 78.0556\%, with widths of 0.571945 and 0.598749 kWh. This closes the scaler-specific comparison on the frozen London cohort, but not site-level calibration. Conformal time-series methods are a candidate extension if their assumptions, calibration cohort, and deployment contract are stated explicitly \cite{stankeviciute2021conformal}.

\subsection{Research-to-serving preprocessing}

The cold-start experiment gives every held-out London household a mean and standard deviation fitted on its own historical training segment; the application instead estimates both statistics from the current 336-hour request window. The preserved artifacts were sufficient to rerun exactly the same cold households, fixed origins, histories, targets, seasonal baselines, calendar features, and checkpoints with the production rolling normalization and output transform. Macro MAE is 0.181945 kWh at 24 hours and 0.194138 kWh at 168 hours under that serving transform, compared with the unchanged research headlines of 0.184928 and 0.197322 kWh. This resolves the identified transform mismatch for the frozen LCL cohort. It does not establish expected accuracy for an arbitrary new site because geography, data quality, time-zone handling, and calibration may differ.

\subsection{Monthly horizon}

The earlier monthly N-BEATS and MultiCycleNet candidates remain invalid for promotion because their cold-start macro $R^2$ values are negative. The replacement Chronos-2 release passes its frozen selection gate with 49.82\% improvement and positive macro $R^2$, then transfers positively to Portugal. Its southern-Morocco diagnostic is worse than the comparator and therefore prevents a broad geographic claim. The repository capability is production-eligible for controlled use, not universally validated.

\subsection{Experiment completeness}

TimePro and TimeMixer++ were attempted but lack completed leaderboard evidence. TSMixer appears in an exploratory weather notebook without a frozen successful metric artifact. MultiCycleNet was adapted locally and requires final primary-reference verification before any claim of exact reproduction. These gaps are reported rather than filled with inferred results.

\section{Product limitations}

The release has verified public web and API URLs, email registration and delivery capability, a healthy email worker, and a production Google identity origin, but that does not make every repository capability operational in the cloud. Simulation, alert, and avatar-cleanup workers are absent; avatars are ephemeral; centralized log retention is not configured; and backup restoration has not been tested. The Azure-generated hostname is not a custom domain, the registry is not yet integrated through workload identity, and the budget is an alerting control rather than an automatic shutdown. Application security headers are implemented and locally verified in the remediated source but require image promotion before they become Azure evidence. Formal accessibility, load, disaster-recovery, sustained provider, and multi-tenant penetration testing remain incomplete. Horizontal multi-replica rate limiting also requires a shared backend rather than process-local counters. The research tree contains historical assets that should remain separated from the deployable product snapshot.

\section{Threats to validity}

\begin{table}[H]
\centering
\caption{Threats to validity and mitigation.}
\label{tab:validity}
\begin{tabularx}{\textwidth}{p{0.19\textwidth}YY}
\toprule
Validity type & Threat & Mitigation and remaining gap \\
\midrule
Internal & Leakage, preprocessing scope, repeated test use & Causal windows, chronological split, train-only transforms, evidence registries; historic notebooks remain separately labelled \\
Research--serve & Research scaler differs from 336-hour serving scaler & Exact frozen-cohort comparison completed; new-site transfer and calibration remain required \\
Construct & Accuracy/classification confused with forecasting value & Continuous target, physical-unit errors, seasonal baselines, household distribution \\
External & London/IHEPC results transferred to Moroccan users & Cold-start household split; local-country evaluation still required \\
Conclusion & Single seed, clock-influenced dynamic sampling, aggregation choice & Frozen seed and common protocol, macro/median/p90/win rate; bitwise retraining and multi-seed finalist study remain desirable \\
Implementation & Third-party repository/version differences & Isolated workers and source notes; incomplete models excluded from rankings \\
Operational & Public endpoint interpreted as complete production maturity & Exact Azure boundary recorded; absent workers, ephemeral storage, monitoring, and recovery tests remain explicit \\
\bottomrule
\end{tabularx}
\end{table}

\section{Discussion}

The project demonstrates a change in what counts as success. At the beginning, success meant matching a 98\% article result. After the audit, success meant constructing a task that could fail honestly. That shift led to lower and more interpretable metrics, stronger baselines, broader household analysis, and explicit artifact limitations.

The final model decisions are credible because they combine causal targets, frozen protocols, strong comparators, distributional metrics, transfer diagnostics, checksummed artifacts, and explicit operational gates. The application adds strict loading and refuses unsupported cadence or history. The decision not to activate the slightly more accurate hourly Chronos challengers is part of this rigor: a production model must improve the whole system contract, not only one metric. These facts support a controlled demonstration and a defensible PFE contribution, but they do not establish universal geographic transfer or site-specific calibration.

The remaining conditional findings also improve rather than weaken the report. They separate verified public availability from full production maturity and prevent a controlled academic deployment from being misrepresented as a globally validated service.
